\chapter{Conclusioni}

In questo elaborato è stata illustrata l'efficacia dell'espansione in multipoli nella gravitazione newtoniana per descrivere in modo rigoroso le deviazioni dalla simmetria sferica, confrontando la decomposizione in armoniche sferiche con il formalismo cartesiano dei tensori STF. L'introduzione dei momenti di multipolo ha fornito degli strumenti analitici per quantificare la deviazione dei corpi celesti dalla simmetria sferica e, applicato ai dati gravimetrici terrestri, questo formalismo permette di quantificare lo schiacciamento rotazionale del pianeta, sperimentalmente dominato dal momento di quadrupolo, e di isolare, nei termini di ordine superiore, le anomalie locali generate dalla complessa geodinamica interna. È stato inoltre approfondito il problema delle deformazioni mareali indotte da sorgenti esterne, dimostrando, sia in ambito classico sia nel limite relativistico, come il campo mareale induca un momento di quadrupolo proporzionale al campo stesso, parametrizzato tramite i numeri di Love gravitazionali, e abbiamo mostrato il peculiare annullamento esatto del coefficiente di deformabilità mareale nel caso dei buchi neri.


I principali punti di forza di questo impianto teorico risiedono nella sua semplicità analitica e nell'ampia generalità applicativa, tuttavia emergono precisi vincoli fisici e operativi: dal punto di vista della modellizzazione, la determinazione dei coefficienti richiede la conoscenza dettagliata della distribuzione di densità e, nel contesto mareale, dell'equazione di stato della materia interna, necessaria per integrare le equazioni di campo all'interno del corpo ed eseguire il corretto raccordo sulla superficie. Inoltre, l'espansione in multipoli mostra evidenti limiti di convergenza quando applicata a configurazioni che si discostano marcatamente dalla simmetria sferica, come le galassie a spirale: in tali scenari, sebbene lo sviluppo resti formalmente valido, la lenta convergenza della serie impone l'inclusione di ordini molto elevati, rendendo il calcolo analitico poco pratico e richiedendo delle necessarie strategie algoritmiche per eliminare le instabilità numeriche connesse al calcolo delle funzioni associate di Legendre.
